\documentclass[10pt,a4paper]{amsart}
\usepackage[margin=19mm]{geometry}
\usepackage{amsmath,amssymb,mathtools}
\usepackage[scr=rsfs,cal=euler]{mathalfa}
\usepackage{xfrac}
\usepackage[unicode,hidelinks,bookmarks=false]{hyperref}
\newcommand{\ie}{i.e.\ }
\newcommand{\cf}{cf.\ }
\newcommand{\resp}{respectively\ }
\newcommand{\Subsection}{\subsection}
\providecommand{\nobreakdash}{\nobreak\hbox{-}}
\setlength{\emergencystretch}{2em}
\multlinegap0pt
\usepackage{amssymb}
\usepackage{mathtools}
\usepackage[shortlabels]{enumitem}
\newtheorem{lemma}{Lemma}
\def\R{{\mathbb{R}}}
\def\pl{{\text{PL}^\mu}}
\title{Your Title Here}
\author{Julien Hermant}
\date{Source: arXiv:2602.10022v2 (CC BY 4.0)}
\begin{document}
\maketitle

\noindent\emph{Source and licence.} Your Title Here. Version arXiv:2602.10022v2 (CC BY 4.0), distributed by arXiv under the Creative Commons Attribution 4.0 International licence, http://creativecommons.org/licenses/by/4.0/, as recorded in the arXiv licence metadata for this exact version and shown on the abstract page https://arxiv.org/abs/2602.10022v2. Full licence notice: \texttt{licenses/arxiv-2602.10022-CC-BY-4.0.txt}.\par\smallskip

\noindent\textit{Editorial scope.} Counted unit: the article's lemma lem:toy\_example\_pl (source0.tex lines 1468--1473). The article's complete original proof of it is delivered with it (source0.tex lines 1474--1587, one \texttt{proof} environment of 114 lines). No mathematics is written, repaired or re-derived: the statement and the proof are the article's own lines, in the article's own notation, and no retained line is substituted or re-flowed.  No further source-local context is needed: every cross-reference of every retained line resolves inside the two delivered spans.  Nothing was omitted from any retained span, so the package carries no omission record.  No retained line contains a citation, so the wrapper carries no bibliography and no bibliography entry.  No retained line contains a figure environment, an image-inclusion command or drawing code, so no image is delivered and the package declares no asset dependency.  Editorial formatting only, in addition: an AMS \texttt{amsart} preamble that restates the article's own declarations and macros used by the retained lines, namely the environments \texttt{lemma} on separate counters, together with the standard packages those lines need.  The retained environments print in this document as Lemma 1 in order of appearance, whereas the article prints the same items with the numbering of its own full document.  The complete native source of the article as distributed in the arXiv e-print for this exact version is bundled unchanged as \texttt{sources/194366/source0.tex}.
\begin{lemma}\label{lem:toy_example_pl}
    For $x,y\in\R,\ \varepsilon>0$ the function $F_{\varepsilon}(x,y)=\tfrac12\,(y-\sin x)^2+\tfrac12\,\varepsilon x^2$ belongs to $\pl$ with
  $$\mu \le \frac{2+\varepsilon
-\sqrt{\varepsilon^2+4}}{2}.$$

\end{lemma}

\begin{proof}
    Consider the function
\[
F_{\varepsilon}(x,y)=\tfrac12\,(y-\sin x)^2+\tfrac12\,\varepsilon x^2,
\qquad x,y\in\R,\ \varepsilon>0.
\]

Let
\[
u(x,y) := y - \sin x.
\]
Then
\[
F_{\varepsilon}=\tfrac12\,(u^2+\varepsilon x^2).
\]

The gradient is
\[
\nabla F_{\varepsilon}=
\begin{pmatrix}
-u\cos x + \varepsilon x\\[3pt]
u
\end{pmatrix},
\qquad
\|\nabla F_{\varepsilon}\|^2 = u^2 + (-u\cos x + \varepsilon x)^2.
\]


Writing both sides as quadratic forms in \(z=(u(x,y),x)^\top\), we obtain
\[
2F_{\varepsilon}
= z^\top 
\begin{pmatrix}
1 & 0\\[3pt]
0 & \varepsilon
\end{pmatrix}z,
\qquad
\|\nabla F_{\varepsilon}\|^2
= z^\top
\begin{pmatrix}
1+\cos^2 x & -\varepsilon\cos x\\[3pt]
-\varepsilon\cos x & \varepsilon^2
\end{pmatrix}z.
\]

Define
\[
N=\begin{pmatrix}1&0\\0&\varepsilon\end{pmatrix},\qquad
M(x)=\begin{pmatrix}
1+\cos(x)^2 & -\varepsilon \cos(x)\\[3pt]
-\varepsilon \cos(x) & \varepsilon^2
\end{pmatrix},
\qquad c:=\cos x\in[-1,1].
\]
So, $f$ belong to $\pl$ with
\begin{equation}
    \begin{aligned}
        \mu = \min_{\substack{z=(u(x,y),x)\neq 0,\\(x,y)\in \R^2}}~ \frac{z^T M(x)z}{z^T N z} &\ge \min_{\substack{z=(z_1,z_2),\\(z_1,z_2)\in \R^2\backslash (0,0)}}~ \frac{z^T M(x)z}{z^T N z} \\&\overset{(i)}{=}  \min_{\substack{z=(z_1,z_2),\\(z_1,z_2)\in \R^2\backslash (0,0)}}~ \frac{z^T N^{-1/2}M(x)N^{-1/2}z}{z^Tz} \\
        &\overset{(ii)}{=} \lambda_{\min}(N^{-1/2}M(x)N^{-1/2}),
    \end{aligned}
\end{equation}
where (i) holds with the reparameterization $z' = N^{-1/2}z$, and (ii) is a property of the Rayleigh coefficient. We consider the matrix
\[
N^{-1/2}M(x)N^{-1/2}
=
\begin{pmatrix}
1+\cos(x)^2 & -\sqrt{\varepsilon}\,\cos(x)\\[3pt]
-\sqrt{\varepsilon}\,\cos(x) & \varepsilon
\end{pmatrix}
:= \begin{pmatrix}
1+c^2 & -\sqrt{\varepsilon}\,c\\[3pt]
-\sqrt{\varepsilon}\,c & \varepsilon
\end{pmatrix}=: B(c),\quad c:= \cos(x) \in [-1,1].
\]
So, we have
$$\lambda_{\min}(N^{-1/2}M(x)N^{-1/2}) \ge \inf_{c\in[-1,1]}\ 
\lambda_{\min}(B(c)).$$

The eigenvalues of \(B(c)\) are
\[
\lambda_\pm(c)
=
\frac{\operatorname{tr}(B(c))
\pm
\sqrt{\operatorname{tr}(B(c))^2 - 4\det(B(c))}}{2}.
\]

We have
\[
\det(B(c))=\varepsilon,
\qquad
\operatorname{tr}(B(c)) = 1 + c^2 + \varepsilon.
\]
Such that we solve
$$\min_{c \in [-1,1]} \frac{1+c^2 + \varepsilon - 
\sqrt{(1+c^2 + \varepsilon)^2 - 4\varepsilon}}{2} =\min_{t \in [0,1]} \frac{1+t + \varepsilon - 
\sqrt{(1+t + \varepsilon)^2 - 4\varepsilon}}{2} := \phi(t).$$
We have for $t \in [0,1]$,
$$2\phi'(t) = 1 - \frac{1+t+\varepsilon}{\sqrt{(1+t + \varepsilon)^2 - 4\varepsilon}} \le 0,$$
such that $\phi$ is decreasing on $[0,1]$. So, the minimum is reached for $t = 1 \Rightarrow c^2 = 1$, which leads to the eigenvalue
$$\inf_{c\in[-1,1]}\ 
\lambda_{\min}(B(c)) = \lambda_{\min}(B(1)) = \frac{2+\varepsilon
-\sqrt{\varepsilon^2+4}}{2}
>0.$$


Therefore, for all \((x,y)\in\R^2\),
\[
\|\nabla F_{\varepsilon}(x,y)\|^2 \;\ge\; 2\mu\, F_{\varepsilon}(x,y),
\]
with 
$$\mu \le \frac{2+\varepsilon
-\sqrt{\varepsilon^2+4}}{2}.$$
\end{proof}

\end{document}
